\chapter{Implementation}

The project is implemented in Python with a modular structure.
The main libraries are NumPy, pandas, SciPy, scikit-learn, matplotlib, and PyTorch.

\section{Code Structure}

\begin{lstlisting}[style=pythonstyle, caption={Main project structure.}]
src/nn_option_pricing/
    black_scholes.py   # analytical formula and payoff
    monte_carlo.py     # efficient Monte Carlo simulation
    dataset.py         # synthetic data generation and storage
    model.py           # PyTorch feed-forward neural network
    training.py        # training loop, scaling, early stopping
    evaluation.py      # regression metrics
    plots.py           # visualizations
    pipeline.py        # end-to-end experiment
scripts/
    run_experiment.py  # command-line entry point
\end{lstlisting}

\section{Computational Efficiency}

The Black-Scholes formula and the dataset generation procedure are implemented in vectorized form with NumPy.
The Monte Carlo method uses the exact terminal distribution of the geometric Brownian motion at time $T$, avoiding unnecessary time discretization.
Moreover, Monte Carlo simulation is batched both over options and over simulated paths, which keeps memory usage under control.

The neural network training uses PyTorch \texttt{DataLoader}, mini-batch gradient descent, input standardization, and target scaling.
These choices make the code suitable both for quick development runs and for larger experiments.

\section{Reproducibility}

To make the experiments reproducible, the project fixes the random seeds of NumPy, Python, and PyTorch.
The artifacts produced by an experiment include:

\begin{itemize}
    \item the generated synthetic dataset;
    \item the model checkpoint;
    \item the input scaler;
    \item the target scaler;
    \item metrics in JSON format;
    \item diagnostic plots.
\end{itemize}
